\chapter{More on Lists and Loops in Python}

\section{Introduction: Real-Life Hook}

Students often deal with multiple collections of items:

\begin{itemize}
    \item A class with multiple students, each having multiple subjects  
    \item Grocery shopping list with items and their quantities  
    \item Scores of students in different tests  
\end{itemize}

Python allows us to \textbf{nest lists} (lists within lists) and \textbf{loop through them efficiently}.

\section{Concept: Nested Lists}

\begin{definition}{Nested List}
A \textbf{nested list} is a list that contains other lists as elements.
\end{definition}

\subsection{Example: Students’ Marks}

\begin{verbatim}
marks = [
    [80, 70, 90],   # Student 1
    [75, 85, 65],   # Student 2
    [90, 95, 100]   # Student 3
]

print("Marks of Student 2:", marks[1])
print("Second subject of Student 3:", marks[2][1])
\end{verbatim}

Output:

\begin{verbatim}
Marks of Student 2: [75, 85, 65]
Second subject of Student 3: 95
\end{verbatim}

- Use \textbf{double indexing} to access elements in nested lists  

\section{Looping Through Nested Lists}

\begin{verbatim}
for i, student in enumerate(marks):
    print("Student", i+1, "marks:")
    for mark in student:
        print(mark, end=" ")
    print()
\end{verbatim}

Output:

\begin{verbatim}
Student 1 marks: 80 70 90
Student 2 marks: 75 85 65
Student 3 marks: 90 95 100
\end{verbatim}

- Outer loop iterates through students  
- Inner loop iterates through marks of each student  

\section{Advanced List Operations}

\subsection{Slicing Lists}

\begin{verbatim}
fruits = ["apple", "banana", "mango", "orange"]
print(fruits[1:3])  # ['banana', 'mango']
print(fruits[:2])   # ['apple', 'banana']
print(fruits[2:])   # ['mango', 'orange']
\end{verbatim}

\subsection{List Functions and Methods}

\begin{verbatim}
numbers = [5, 2, 9, 1]
print("Max:", max(numbers))
print("Min:", min(numbers))
print("Sum:", sum(numbers))
numbers.sort()
print("Sorted:", numbers)
\end{verbatim}

\section{Mini-Project: Student Scores Table}

Create a program to store marks of 3 students in 3 subjects and print total marks:

\begin{verbatim}
students = [
    [80, 70, 90],
    [75, 85, 65],
    [90, 95, 100]
]

for i, student in enumerate(students):
    total = sum(student)
    print("Total marks of Student", i+1, "=", total)
\end{verbatim}

Output:

\begin{verbatim}
Total marks of Student 1 = 240
Total marks of Student 2 = 225
Total marks of Student 3 = 285
\end{verbatim}

\section{Step-by-Step Diagram: Nested Lists Loop Flow}

\begin{center}
\begin{tikzpicture}[node distance=1.5cm]
\node (start) [draw, ellipse, fill=green!20] {Start Program};
\node (outer) [draw, rectangle, fill=yellow!20, below of=start] {Loop Through Outer List};
\node (inner) [draw, rectangle, fill=orange!20, below of=outer] {Loop Through Inner List};
\node (process) [draw, rectangle, fill=orange!20, below of=inner] {Process Each Element};
\node (end) [draw, ellipse, fill=green!20, below of=process] {End Program};
\draw[->] (start) -- (outer);
\draw[->] (outer) -- (inner);
\draw[->] (inner) -- (process);
\draw[->] (process) -- (end);
\end{tikzpicture}
\end{center}

\section{Key Concepts}

\begin{itemize}
    \item Nested lists store lists inside lists for hierarchical data  
    \item Double indexing accesses elements in nested lists  
    \item Loops can be nested to iterate over complex lists  
    \item Advanced operations: slicing, max, min, sum, sort  
\end{itemize}

\section{End-of-Chapter Exercises}

\begin{exercise}{Very Short Questions}
\begin{enumerate}
    \item What is a nested list?  
    \item How do you access the first element of a nested list?  
    \item Name two list methods to manipulate numbers.  
\end{enumerate}
\end{exercise}

\begin{exercise}{Short Answer Questions}
\begin{enumerate}
    \item Write a Python program to print the second subject marks of each student.  
    \item Explain why nested loops are needed for nested lists.  
    \item Give an example of a nested list in daily life.  
\end{enumerate}
\end{exercise}

\begin{exercise}{Activity-Based Questions}
\begin{enumerate}
    \item Store marks of 4 students in 3 subjects each and print the average marks of each student.  
    \item Create a list of lists representing a shopping list (item name and quantity) and print them.  
    \item Make a program that sorts each student’s marks in ascending order using a loop.  
\end{enumerate}
\end{exercise}

\section{Teacher Notes}

\marginnote{
\textbf{Learning Objectives:}  
\begin{itemize}
    \item Students understand nested lists and double indexing  
    \item Students can use nested loops to process complex lists  
    \item Students can apply advanced list operations (slicing, max, min, sum, sort)  
\end{itemize}

\textbf{Teaching Tips:}  
\begin{itemize}
    \item Use real-life examples: students’ marks, shopping lists  
    \item Demonstrate nested loops step-by-step  
\end{itemize}

\textbf{Assessment:}  
\begin{itemize}
    \item Check correctness of nested list loops  
    \item Observe understanding of list operations and slicing  
\end{itemize}
}